\chapter{电子成像、收集与参数反演}\label{chap:electron-imaging-inversion}

源与材料核已经给出理想损失及辐射概率。真实图像还包含收集角、能量分辨率和计数噪声；这些过程进入同一个前向模型以后，才能判断反演参数是否唯一。

\section{成像、反演与收集通道}

前向理论给出的场或末态分布还不是实验数据。实验真正记录的是像素、能量 bin、到达时间或计数事件，因此理论必须再经过收集光学和探测器。把一组物理参数 $\boldsymbol\vartheta$ 映射到数据概率分布 $p(y|\boldsymbol\vartheta)$ 的计算链称为\emph{前向模型}；给定实际数据 $y$，反过来约束或估计 $\boldsymbol\vartheta$ 的问题称为\emph{反演}。反演是否唯一取决于不同参数是否会在前向模型中留下可区分的数据变化，而不只取决于方程能否数值拟合。

有限数值孔径只收集部分辐射通道，谱仪把真实能量与仪器响应卷积，计数噪声再把连续概率分布变成随机样本。若探测器具有空间成像功能，同一思想中的空间条件响应通常称为点扩散函数（point-spread function, PSF）：理想点源经过有限孔径和像差后在探测面形成的归一化强度核；在空间平移不变的近似下，实物像就是理想像与 PSF 的卷积。能谱仪的 $r_E(E_d|E_f)$ 是完全类似的“能量轴 PSF”。因此成像与参数反演必须把“电子源—材料响应—传播—探测”串成同一个概率模型，才能把真实物理响应与仪器模糊区分开。

\subsection{近场能谱成像}\label{subsec:nearfield-spectrum-imaging-dp19}

扫描式 PINEM 在每个横向位置记录一条电子能谱。用 $(x,y)$ 表示扫描坐标，$E_d$ 表示仪器读出的能量，整组记录便写为 $I(x,y,E_d)$。求这组数据时，先给定入射能量分布，再计算该位置的光学近场如何调制电子，最后加入谱仪对能量的响应。这样便可判断观测到的谱宽来自入射源、近场作用还是仪器分辨率。

扫描坐标记为
\begin{equation*}
\mathbf R=(x,y),
\qquad [\mathbf R]={\rm m}.
\end{equation*}
对中心速度 $v_0$ 的电子和单频近场，定义位置依赖的无量纲 PINEM 耦合
\begin{equation}
\boxed{
\beta(\mathbf R)
=\frac{q_{\mathrm e}}{\hbar\omega}
\int_{-\infty}^{\infty}\dd z\,
E_{\omega,z}(\mathbf R,z)
\ee^{-\ii\omega z/v_0}.}
\label{eq:pinem-spatial-beta-dp19}
\end{equation}
这里 $E_{\omega,z}$ 是同步纵向复电场幅度，单位 ${\rm V\,m^{-1}}$。积分中的 $q_{\mathrm e} E_{\omega,z}\dd z$ 是焦耳，除以单光子能量 $\hbar\omega$ 后 $\beta$ 无量纲。指数因子只保留与电子飞行相位 $\omega z/v_0$ 同步的纵向 Fourier 分量，因此 $\beta(\mathbf R)$ 同时携带局域场幅度和相位信息。

在低反冲、双边能量平移模型中，固定 $\mathbf R$ 时的理想边带条件概率为
\begin{equation*}
P(\ell|\mathbf R)
=J_\ell^2(2|\beta(\mathbf R)|),
\qquad
\sum_{\ell\in\mathbb Z}P(\ell|\mathbf R)=1.
\end{equation*}
$\ell\in\mathbb Z$ 是净交换光量子数；真实末态能量相对入射能量移动 $\ell\hbar\omega$。

电子源具有有限初始能散。定义归一化入射能量概率密度
\begin{equation*}
F_0(E)\ge0,
\qquad
\int\dd E\,F_0(E)=1,
\qquad
[F_0]={\rm J^{-1}}.
\end{equation*}
谱仪对真实末态能量 $E_f$ 的读出分成两个物理问题：电子是否被记录，以及被记录后读成哪个能量。前者由探测效率控制，后者由条件线形控制。令
\begin{equation*}
r_E(E_d|E_f)\ge0,
\qquad
\int\dd E_d\,r_E(E_d|E_f)=1,
\qquad
[r_E]={\rm J^{-1}},
\end{equation*}
以及无量纲效率
\begin{equation*}
0\le\mathcal Q_E^{\rm det}(E_f)\le1.
\end{equation*}
于是有量纲点击响应密度为
\begin{equation*}
R_E(E_d|E_f)
\equiv\mathcal Q_E^{\rm det}(E_f)r_E(E_d|E_f),
\qquad
\int\dd E_d\,R_E(E_d|E_f)=\mathcal Q_E^{\rm det}(E_f).
\end{equation*}

在入射电子可由经典能量混合 $F_0(E)$ 描述时，扫描位置 $\mathbf R$ 处的\emph{点击能谱密度}为
\begin{equation}
\boxed{
 p_{\rm click}(E_d|\mathbf R)
 =\sum_{\ell\in\mathbb Z}\int\dd E\,
 R_E(E_d|E+\ell\hbar\omega)
 J_\ell^2\!\left(2|\beta(\mathbf R)|\right)
 F_0(E).}
\label{eq:pinem-spatial-detector-forward-dp19}
\end{equation}
$[p_{\rm click}]={\rm J^{-1}}$。这个卷积可以直接按条件概率读取：固定真实入射能量 $E$ 与扫描位置 $\mathbf R$ 时，边带量子数和探测读出的联合密度为
\begin{equation*}
p_{\rm click}(E_d,\ell|E,\mathbf R)
=R_E(E_d|E+\ell\hbar\omega)
J_\ell^2\!\left(2|\beta(\mathbf R)|\right).
\end{equation*}
对互斥边带 $\ell$ 求和、再按源分布 $F_0(E)$ 平均，正好得到\eqnrefs{eq:pinem-spatial-detector-forward-dp19}；这一步就是条件概率的边缘化。其积分
\begin{equation*}
P_{\rm click}(\mathbf R)
\equiv\int\dd E_d\,p_{\rm click}(E_d|\mathbf R)
\le1
\end{equation*}
是该扫描点一次电子被记录的总概率。若实验数据已经条件在“发生点击”这一事件上，则用于形状拟合的归一化谱为
\begin{equation*}
p(E_d|\mathbf R,{\rm click})
=\frac{p_{\rm click}(E_d|\mathbf R)}{P_{\rm click}(\mathbf R)}.
\end{equation*}

\eqnrefs{eq:pinem-spatial-detector-forward-dp19} 把三个实验层次分离：$F_0$ 属于电子源，$\beta(\mathbf R)$ 属于相互作用区，$R_E$ 属于探测器。若入射电子保留能量相干，则 $F_0(E)$ 只给出密度矩阵对角元，不能描述完整初态；此时必须使用量子仪器\eqnrefs{eq:general_indirect_measurement_instrument}，让不同能量路径在振幅层进入同一个探测结果。

若结构还提供可分辨横向动量转移，末态应标记为 $(\ell,\mathbf p_\perp)$。定义联合条件概率密度
\begin{equation*}
P_\ell(\mathbf p_\perp|\mathbf R)
=|\mathcal A_\ell(\mathbf p_\perp;\mathbf R)|^2,
\end{equation*}
其中
$[P_\ell(\mathbf p_\perp|\mathbf R)]={\rm momentum^{-2}}$，所以 $P_\ell\dd^2 p_\perp$ 无量纲。能量轴读取净交换量子数；横向动量轴读取结构承担的平行动量差。二维谱因而同时解析局域同步 Fourier 分量以及横向平移对称性破缺所产生的动量转移。\cite{Vanacore2018Attosecond}

\begin{figure}[htbp]
\centering
\includegraphics[width=0.68\linewidth]{pinem_experiment_geometry.png}
\caption{PINEM 的基本实验几何。扫描电子在纳米结构附近采样位置依赖的同步纵向近场，产生离散能量平移；入射能散和谱仪响应随后把理想边带映射成\eqnrefs{eq:pinem-spatial-detector-forward-dp19} 的点击能谱。}
\label{fig:pinem-experiment-interface-dp19}
\end{figure}

\subsection{实验尺度与反演}\label{subsec:experimental-scales-inversion-dp19}

拟合得到一组参数，并不一定说明数据能够唯一确定它们。例如，源能宽增大与仪器分辨率变差都可能使谱线展宽；如果它们对所有读数的作用完全相同，就无法仅凭这组数据区分二者。将各预测读数对参数求偏导，组成前向模型的 Jacobian 矩阵；不同参数对应的列是否独立，可以判断局部上能否区分这些参数。若还把各读数的计数波动计入，就得到第十章定义的 Fisher 信息，用来估计可辨识参数的统计精度。

模型选择由不同物理层的无量纲尺度共同决定。电子运动学决定中心速度和 de Broglie 波长；有限长度相位失配决定相互作用能否相干积累；反冲非谐性决定等间隔能量格是否可用；探测器分辨率决定相邻边带能否区分。把这些尺度分开可以避免用单一的“电子能量高”或“耦合强”替代真正的近似条件。

设电子中心动能为
\begin{equation*}
K_0=E_0-m_{\mathrm e} c^2,
\qquad [K_0]={\rm J}.
\end{equation*}
由质量壳关系
\begin{equation*}
\gamma_0=1+\frac{K_0}{m_{\mathrm e} c^2},
\qquad
v_0=c\sqrt{1-\gamma_0^{-2}},
\qquad
p_0=\gamma_0m_{\mathrm e} v_0
\end{equation*}
得到 de Broglie 波长
\begin{equation*}
\lambda_{\rm dB}=\frac{2\pi\hbar}{p_0}.
\end{equation*}
\tabref{tab:electron_kinematic_scales} 列出常见动能对应的运动学尺度。

电子源的能量标准差记为 $\sigma_{E,{\rm in}}$，单位焦耳。为了比较能量带宽与时间尺度，先定义由带宽构造的特征时间
\begin{equation*}
\tau_E\equiv\frac{\hbar}{\sigma_{E,{\rm in}}}.
\end{equation*}
$\tau_E$ 本身只是一个严格定义的带宽时间尺度；具体的一阶相干时间还取决于谱线形与所采用的相干宽度定义。例如变换极限高斯波包满足相应的高斯时频宽度乘积，不能把任意谱形的相干时间一律等同于 $\tau_E$。
逐发电子--光到达时间抖动记为 $\sigma_t$，单位秒；它产生光学相位噪声尺度 $\omega\sigma_t$。$\sigma_{E,{\rm in}}$ 描述单电子或电子系综的谱宽，而 $\sigma_t$ 描述不同实验重复之间的参考时钟不确定性，两者进入密度矩阵的方式不同。

\begin{table}[H]
\centering
\caption{若干常用电子动能对应的相对论运动学尺度。}
\label{tab:electron_kinematic_scales}
\begin{tabular}{rrrr}
\toprule
$K_0$ (keV) & $\gamma_0$ & $v_0/c$ & $\lambda_{\mathrm dB}$ (pm)\\
\midrule
10 & 1.0196 & 0.1950 & 12.20\\
30 & 1.0587 & 0.3284 & 6.98\\
100 & 1.1957 & 0.5482 & 3.70\\
200 & 1.3914 & 0.6953 & 2.51\\
300 & 1.5871 & 0.7765 & 1.97\\
\bottomrule
\end{tabular}
\end{table}

对有限长度 $L$ 的相互作用区，设纵向相位失配为 $\Delta\kappa$，单位 ${\rm m^{-1}}$；定义
\begin{equation*}
\epsilon_{\rm sync}\equiv|\Delta\kappa|L.
\end{equation*}
$\epsilon_{\rm sync}\ll1$ 时整个长度内相位基本同向积累，$\epsilon_{\rm sync}\gtrsim1$ 时不同位置贡献开始相消。

对本节纵向 PINEM 交换，反冲频率是 $\omega_{\rm r}=\hbar\kappa_z^2/(2m_\parallel)$；驻波横向衍射另用其几何对应的 $\omega_{\rm KD}$，不能直接替换。若相互作用持续时间为 $T$，相邻交换边的反冲非谐性可用
\begin{equation*}
\epsilon_{\rm rec}\equiv2\omega_{\rm r}T
\end{equation*}
衡量；当该相位达到一阶量级时，平移对称等间隔格不能再忽略相邻边的不同失谐。

定义源能散和仪器分辨率合成后的有效能量宽度 $\sigma_E^{({\rm eff})}$，单位焦耳。相邻 PINEM 边带间隔为 $\hbar\omega$，因此谱分辨判据为
\begin{equation*}
\epsilon_{\rm det}
\equiv\frac{\sigma_E^{({\rm eff})}}{\hbar\omega}.
\end{equation*}
$\epsilon_{\rm det}\ll1$ 表示相邻边带在能量轴上可分辨；$\epsilon_{\rm det}\gtrsim1$ 时单个边带结构被卷积成连续展宽，但完整前向模型仍然可以包含这些未分辨贡献。

若只使用总能谱二阶矩估计 $|\beta|$，还需要更强的统计假设。设入射能量、PINEM 边带阶数和谱仪加性读出噪声彼此独立，谱仪噪声零均值；定义它们的标准差分别为 $\sigma_{E,{\rm in}}$、$\sqrt{\operatorname{Var}(\ell)}\hbar\omega$ 与 $\sigma_{E,{\rm inst}}$。在纯相位低反冲 PINEM 中
\begin{equation*}
\operatorname{Var}(\ell)=2|\beta|^2,
\end{equation*}
因此
\begin{equation}
\boxed{
\sigma_{E,{\rm meas}}^2
=\sigma_{E,{\rm in}}^2
+\sigma_{E,{\rm inst}}^2
+2(\hbar\omega)^2|\beta|^2.}
\label{eq:pinem_energy_variance_addition}
\end{equation}
反解得到
\begin{equation}
\boxed{
|\beta|
=
\frac{\sqrt{
\sigma_{E,{\rm meas}}^2
-\sigma_{E,{\rm in}}^2
-\sigma_{E,{\rm inst}}^2}}
{\sqrt2\,\hbar\omega}.}
\label{eq:beta_from_variance_experiment}
\end{equation}
根号内必须非负；若实验估计因有限样本出现负值，说明二阶矩误差已大于可辨识 PINEM 展宽，不能把该样本代入实数 $|\beta|$ 公式。

方差估计只使用一个统计量，因而会丢失完整边带线形。定量反演应使用\eqnrefs{eq:pinem-spatial-detector-forward-dp19} 或一般实验前向模型~\eqnrefs{eq:general_experimental_forward_model}。若 Fisher 信息矩阵在某参数方向上有零模，重复同一设置不会增加该方向的信息；改变扫描位置、参考相位、传播距离或耦合强度才会改变前向模型的导数空间。

\subsection{有限分辨率、有限窗口与近场复振幅反演}
\label{sec:electron-detector-inversion-worked}

能谱仪记录的是每个能量区间内的电子数。要由理想 PINEM 边带计算这些计数，需要先加入入射能宽和仪器误差，再对各区间积分。以\eqnrefs{eq:pinem-spatial-detector-forward-dp19}为起点，设入射能量偏差与仪器加性误差为彼此独立的零均值高斯变量，标准差分别是 $\sigma_{E,\rm in}$、$\sigma_{E,\rm inst}$。记 $\varepsilon_d=E_d-E_0$ 为相对于源中心能量的读出偏差，并假设整个建模能区中的探测效率为常数 $\mathcal Q_E^{\rm det}$。在已经探测到电子这一条件下，读出能量的概率密度为
\begin{equation}
 p(\varepsilon_d|{\rm click})
 =\sum_{\ell\in\mathbb Z}J_\ell^2(2|\beta|)
 \frac{\exp[-(\varepsilon_d-\ell\hbar\omega)^2/(2\sigma_E^2)]}{\sqrt{2\pi}\sigma_E},
 \qquad \sigma_E^2=\sigma_{E,\rm in}^2+\sigma_{E,\rm inst}^2.
 \label{eq:electron-detector-gaussian-comb}
\end{equation}
这一步还要求源能宽内的 $\beta$ 变化可忽略，并采用经典低反冲 PINEM；如果能量改变会显著改变速度、同步或耦合，应回到含入射能量依赖的积分。$p$ 的量纲为能量的负一次方，不能把它在峰顶的高度当作该边带概率。

对第 $j$ 个读出窗口 $[a_j,b_j]$，先定义标准正态分布函数
$\Phi_N(x)=(2\pi)^{-1/2}\int_{-\infty}^x\ee^{-u^2/2}\dd u$，再积分得到无量纲概率
\begin{equation}
 p_j=\sum_\ell J_\ell^2(2|\beta|)
 \left[\Phi_N\!\left(\frac{b_j-\ell\hbar\omega}{\sigma_E}\right)
       -\Phi_N\!\left(\frac{a_j-\ell\hbar\omega}{\sigma_E}\right)\right].
 \label{eq:electron-detector-bin-probabilities}
\end{equation}
若有 $N_{\rm in}$ 个入射电子、背景已扣除，窗口的期望计数为 $N_{\rm in}\mathcal Q_E^{\rm det}p_j$。若窗口未覆盖整个实轴，还应显式保留“点击但在窗外”的概率 $1-\sum_jp_j$；它不同于未点击概率 $1-\mathcal Q_E^{\rm det}$。

\begin{exbox}[由允许误分率解析确定分辨率要求]
设相邻边带间隔为 $\hbar\omega$、有效高斯标准差为 $\sigma_E$，以两中心的中点划分窗口。求给定真实边带的电子被分到其他窗口的概率。

将中心平移到零，正确窗口为 $[-\hbar\omega/2,\hbar\omega/2]$。左右高斯尾相等，所以
\begin{equation*}
 P_{\rm mis}=2\Phi_N\!\left(-\frac{\hbar\omega}{2\sigma_E}\right).
\end{equation*}
给定容许值 $0<\varepsilon<1$，利用 $\Phi_N$ 单调递增并反解，得到
\begin{equation*}
 P_{\rm mis}\le\varepsilon
 \ \Longleftrightarrow\ 
 \sigma_E\le\frac{\hbar\omega}{2\Phi_N^{-1}(1-\varepsilon/2)}.
\end{equation*}
这里 $\Phi_N^{-1}(y)$ 是满足 $\Phi_N(x)=y$ 的分位点。若希望避免求反函数，还可由正态尾上界
$\Phi_N(-x)<\ee^{-x^2/2}/(\sqrt{2\pi}x)$（$x>0$）得到
\begin{equation*}
 P_{\rm mis}<\sqrt{\frac8\pi}\frac{\sigma_E}{\hbar\omega}
 \exp\!\left[-\frac{(\hbar\omega)^2}{8\sigma_E^2}\right].
\end{equation*}
尾上界来自在积分中用 $u/x\ge1$ 替代一，再积分 $u\ee^{-u^2/2}$。它显示良好分辨区的误分率随间隔宽度比呈指数下降。$P_{\rm mis}$ 是给定真实边带的条件误差；实际某个窗口的占据还须叠加所有边带的输入权重。
\end{exbox}

完整能窗内，高斯卷积只增加一个独立方差，因此
\begin{equation*}
 \sigma_{E,\rm meas}^2=\sigma_E^2+2(\hbar\omega)^2|\beta|^2.
\end{equation*}
若只保留有限读出窗口 $W$，数据却被条件在 $\varepsilon_d\in W$ 上，此时求出的矩是
\begin{equation}
 \langle\varepsilon_d^k\rangle_W
 =\frac{\int_W\varepsilon_d^k p(\varepsilon_d)\dd\varepsilon_d}
        {\int_Wp(\varepsilon_d)\dd\varepsilon_d}.
 \label{eq:electron-detector-windowed-moment}
\end{equation}
分母与 $\beta$ 有关，而且窗外的高阶边带被更强的 $\varepsilon_d^2$ 权重放大；所以少量丢失概率也可能带来明显方差偏差。应按实际窗口计算前向概率，而不是直接把条件方差代入完整谱的公式。

\figref{fig:electron-detector-forward} 显示仪器展宽和窗口截断如何改变同一个物理谱。
\begin{figure}[htbp]
 \centering
 \includegraphics[width=0.92\linewidth]{electron_detector_forward.pdf}
 \caption{同一 $|\beta|=1.20$ 的电子谱经过不同仪器分辨率与有限能窗。左图保留相同真实边带权重，只改变高斯有效宽度；右图用完整方差公式处理逐渐扩大的截断窗，显示朴素反演的偏差。带宽模糊不等于电子相干已经消失。}
 \label{fig:electron-detector-forward}
\end{figure}

即使分辨率无限好，单次经典 PINEM 能谱也只依赖 $|\beta|$，不能给出复近场相位。要读出相位，应增加一个已标定、与待测场相干的参考。若两次作用之间的逐阶反冲相位可忽略，双边平移生成元对易，故
\begin{equation}
 S_{\rm tr}(\beta_R\ee^{\ii\varphi_R})S_{\rm tr}(\beta_s)
 =S_{\rm tr}(\beta_s+\beta_R\ee^{\ii\varphi_R}),\qquad \beta_R>0.
 \label{eq:electron-nearfield-reference-composition}
\end{equation}
这里 $\beta_s$ 是待测的复耦合，$\beta_R$ 是以参考相位原点定义的实正幅度。中间有不可忽略漂移时，应保留 $S_RU_{\rm drift}S_s$ 的顺序；不能继续把两个耦合简单相加。

在完整能窗与已知基线方差下，每个参考相位的能谱给出
\begin{equation}
 q_R(\varphi_R)=\frac{\sigma_{E,\rm meas}^2(\varphi_R)-\sigma_E^2}{2(\hbar\omega)^2}
 =|\beta_s+\beta_R\ee^{\ii\varphi_R}|^2.
 \label{eq:electron-nearfield-quadrature-data}
\end{equation}
$q_R$ 无量纲，是由测得方差构造的数据量，不表示电子电荷。四相位作差给
\begin{equation}
 \Re\beta_s=\frac{q_R(0)-q_R(\pi)}{4\beta_R},\qquad
 \Im\beta_s=\frac{q_R(\pi/2)-q_R(3\pi/2)}{4\beta_R}.
 \label{eq:electron-nearfield-four-phase-field}
\end{equation}
实部和虚部又称为相对于参考的两个\emph{正交分量}；改变共同相位原点会旋转这两个坐标，实验真正确定的是待测场与参考的相对相位。

\begin{lemma}[四参考相位消除强度背景]
设已校准参考幅 $\beta_R>0$，未知复幅 $\beta_s=x_s+\ii y_s$，测量量为 $q_R(\varphi)=|\beta_s+\beta_Re^{\ii\varphi}|^2$，背景在四次测量中相同。相差 $\pi$ 的差分分别给 $4\beta_Rx_s$ 与 $4\beta_Ry_s$，从而\eqnrefs{eq:electron-nearfield-four-phase-field}重建复幅。参考幅为零或背景随扫描变化时不适用该直接反演。
\end{lemma}
\begin{proof}
写 $\beta_s=x_s+\ii y_s$，则
\begin{equation*}
 q_R(\varphi_R)=x_s^2+y_s^2+\beta_R^2
 +2\beta_R(x_s\cos\varphi_R+y_s\sin\varphi_R).
\end{equation*}
相差 $\pi$ 的设置具有相同的 $x_s^2+y_s^2+\beta_R^2$，差值只留下线性的某一个正交分量，因而得到\eqnrefs{eq:electron-nearfield-four-phase-field}。这与前节重建电子 $\rho_{01}$ 的四相位算法形式相似，但未知对象不同：这里重建的是\emph{作用区的复场积分}，前节重建的是\emph{入射电子的密度矩阵元}，不能混用二者的校准系数。
\end{proof}


\begin{exbox}[复近场反演的闭式自检与误差传播]
令 $q_0,q_{\pi/2},q_\pi,q_{3\pi/2}$ 为四相位方差数据，参考幅度 $\beta_R>0$ 已知，求复耦合并检查四组数据能否来自同一个稳定场。

写 $\beta_s=x_s+\ii y_s$，令 $A_s=x_s^2+y_s^2+\beta_R^2$，则四组理想数据依次为
\begin{equation*}
 A_s+2\beta_Rx_s,\quad A_s+2\beta_Ry_s,\quad
 A_s-2\beta_Rx_s,\quad A_s-2\beta_Ry_s.
\end{equation*}
作差给出\eqnrefs{eq:electron-nearfield-four-phase-field}，作和则提供两类自检：
\begin{align*}
 q_0+q_\pi&=q_{\pi/2}+q_{3\pi/2},\\
 \frac{(q_0-q_\pi)^2+(q_{\pi/2}-q_{3\pi/2})^2}{16\beta_R^2}
 &=\frac{q_0+q_{\pi/2}+q_\pi+q_{3\pi/2}}4-\beta_R^2.
\end{align*}
第一式检查相位无关背景是否稳定，第二式检查由干涉差值恢复的 $|\beta_s|^2$ 是否与平均强度一致。其成立条件包括稳定相位、正确基线方差及有效的完整能窗。

若四组 $q$ 的估计误差独立且同方差 $\sigma_q^2$，参考本身视为已知，则线性作差给
\begin{equation*}
 \operatorname{Var}\widehat x_s=\operatorname{Var}\widehat y_s
 =\frac{\sigma_q^2}{8\beta_R^2},\qquad
 \operatorname{Cov}(\widehat x_s,\widehat y_s)=0.
\end{equation*}
这里固定的是 $q$ 的误差，不能据此认为参考越强越好：真实计数下 $\sigma_q$、可收集能窗和参考标定误差都可能随参考幅度变化。参考相位的共同偏移只旋转重建的复平面坐标，所确定的是相对相位。
\end{exbox}

恢复 $\beta_s(\mathbf R)$ 仍只恢复了\eqnrefs{eq:pinem-spatial-beta-dp19} 中沿轨迹的同步场积分，而不是每个 $z$ 处的 $E_{\omega,z}(\mathbf R,z)$。任何附加场 $h(z)$ 只要满足 $\int h(z)\ee^{-\ii\omega z/v_0}\dd z=0$，就不会改变同一速度下的 $\beta_s$。因此从二维扫描重建完整三维近场，需要额外几何模型、速度或角度扫描；提高同一个设置的计数不能移除这类零空间。



\subsection{相位读出与收集通道}\label{subsec:phase-collection-dp19}

前面利用参考调制读取近场的复相位。另一种常见装置让两条电子路径重新相遇，从路径干涉中读取相位。对于辐射光，还需要计算有限孔径实际收集了多少。这两种读出都可以用前章的测量算符处理：先求干涉端口的概率，再求光子进入收集模式的概率。

\begin{proposition}[二路径偏迹、干涉可见度与参考相位读出]
先考虑最小电子路径干涉仪。路径 Hilbert 空间由正交基 $|0\rangle_p,|1\rangle_p$ 张成。参考臂对应归一化环境态 $|\xi_0\rangle$，作用臂经过样品后对应 $|\xi_1(\varphi)\rangle$，其中 $\varphi$ 是待估复近场参数所产生的相位。联合态取
\begin{equation}
\boxed{
|\Psi\rangle
=\frac{|0\rangle_p|\xi_0\rangle
+|1\rangle_p|\xi_1(\varphi)\rangle}{\sqrt2}.}
\label{eq:path_environment_joint_state_dp68}
\end{equation}
在
\[
|\pm\rangle_p
=\frac{|0\rangle_p\pm|1\rangle_p}{\sqrt2}
\]
基底读出，得到
\begin{equation}
\boxed{
P_\pm
=\frac12\left[
1\pm\operatorname{Re}
\langle\xi_0|\xi_1(\varphi)\rangle
\right].}
\label{eq:path-phase-readout-dp19}
\end{equation}
重叠
\[
\mathcal V\ee^{\ii\varphi_{01}}
\equiv\langle\xi_0|\xi_1\rangle,
\qquad 0\le\mathcal V\le1,
\]
该干涉项同时编码可见度 $\mathcal V$ 与相对相位 $\varphi_{01}$。若参考臂再扫描已知相位 $\varphi_R$，令作用臂相对于参考臂额外获得因子 $\ee^{\ii\varphi_R}$，则干涉项严格变为
\[
\operatorname{Re}\!\left[\ee^{\ii\varphi_R}\langle\xi_0|\xi_1\rangle\right]
=\mathcal V\cos(\varphi_{01}+\varphi_R).
\]
因此，取相差 $\pi/2$ 的两个参考相位即可分别读取复重叠的两个正交象限。



任意二路径密度矩阵满足 $p_0+p_1=1$、$|\rho_{01}|^2\le p_0p_1$。采用路径一的参考相位 $e^{+\ii\varphi_R}$，一般无损端口读出由本证明的 Born 二次型给出；对称输入及 $50{:}50$ 端口时，$P_+-P_-=\mathcal V\cos(\varphi_{01}+\varphi_R)$，即\eqnrefs{eq:path-phase-readout-dp19}。混态环境亦由相应偏迹相干元决定。
\end{proposition}
\begin{proof}
先不假定两条路径等幅，也不假定环境末态是纯态。二路径自由度的最一般约化密度矩阵可写成
\[
\hat\rho_p
=
\begin{pmatrix}
p_0 & \rho_{01}\\
\rho_{01}^* & p_1
\end{pmatrix},
\qquad
p_0,p_1\ge0,
\qquad
p_0+p_1=1,
\qquad
|\rho_{01}|^2\le p_0p_1.
\]
其中非对角元 $\rho_{01}$ 同时包含路径间的初始相干以及环境对两条路径的可区分性。若联合态来自
\[
|\Psi\rangle
=
c_0|0\rangle|\xi_0\rangle
+
c_1|1\rangle|\xi_1\rangle,
\]
对环境取偏迹得到
\[
p_0=|c_0|^2,
\qquad
p_1=|c_1|^2,
\qquad
\rho_{01}
=
c_0c_1^*
\langle\xi_1|\xi_0\rangle.
\]
因此环境态越可区分，$|\langle\xi_1|\xi_0\rangle|$ 越小，路径相干就越弱。

在读出前，给路径 $|1\rangle$ 加一个可控参考相位 $+\varphi_R$，
\[
\hat U_R
=
|0\rangle\langle0|
+
\ee^{+\ii\varphi_R}|1\rangle\langle1|.
\]
相位扫描后的密度矩阵为
\[
\hat\rho_p(\varphi_R)
=
\hat U_R\hat\rho_p\hat U_R^\dagger
=
\begin{pmatrix}
p_0 & \rho_{01}\ee^{-\ii\varphi_R}\\
\rho_{01}^*\ee^{+\ii\varphi_R} & p_1
\end{pmatrix}.
\]

再通过一个一般无损二端口。若读出态写成
\[
|+\rangle
=
\cos\frac\theta2\,|0\rangle
+
\ee^{+\ii\chi}
\sin\frac\theta2\,|1\rangle,
\]
\[
|-\rangle
=
-\ee^{-\ii\chi}
\sin\frac\theta2\,|0\rangle
+
\cos\frac\theta2\,|1\rangle,
\]
则 Born 概率为
\begin{align*}
P_+
&=
\langle+|\hat\rho_p(\varphi_R)|+\rangle\\
&=
p_0\cos^2\frac\theta2
+
p_1\sin^2\frac\theta2
+
\sin\theta\,
\operatorname{Re}
\!\left[
\rho_{01}
\ee^{+\ii(\chi-\varphi_R)}
\right],
\end{align*}
以及
\[
P_-=1-P_+.
\]
这个式子已经显示出三个不同实验效应：路径人口不平衡由 $p_0-p_1$ 控制，分束比由 $\theta$ 控制，相干干涉则只由复数 $\rho_{01}$ 决定。

正文采用的是对称输入和 $50{:}50$ 读出，
\[
c_0=c_1=\frac1{\sqrt2},
\qquad
\theta=\frac\pi2,
\qquad
\chi=0.
\]
此时
\[
\rho_{01}
=
\frac12\langle\xi_1|\xi_0\rangle,
\]
从而
\[
P_\pm
=
\frac12
\left[
1
\pm
\operatorname{Re}
\!\left(
\ee^{+\ii\varphi_R}
\langle\xi_0|\xi_1\rangle
\right)
\right].
\]
若暂时取 $\varphi_R=0$，便得到正文
\[
\boxed{
P_\pm
=
\frac12
\left[
1\pm
\operatorname{Re}
\langle\xi_0|\xi_1\rangle
\right].
}
\]
进一步写
\[
\langle\xi_0|\xi_1\rangle
=
\mathcal V\ee^{+\ii\varphi_{01}},
\qquad
0\le\mathcal V\le1,
\]
则扫描参考相位给
\[
P_+-P_-
=
\mathcal V
\cos(\varphi_{01}+\varphi_R).
\]


在对称纯态的特例中，也可直接验证：路径投影后的环境振幅为 $(|\xi_0\rangle\pm|\xi_1\rangle)/2$。取范数平方即得两个端口的概率；Cauchy--Schwarz 给 $|\langle\xi_0|\xi_1\rangle|\le1$。插入参考相位后，重叠乘以已知单位相位，所以两次正交设置给出其实部和虚部。
\end{proof}

\begin{definition}[条件单激发的收集投影]
再考虑电子后选择预示的单光子是否进入实际光路。设后选择后的归一化总单激发态为 $|\psi_1\rangle$。定义由收集数值孔径、偏振选择、频率滤波和空间模式匹配共同确定的正交投影算符
\begin{equation*}
\Pi_{\rm col},
\qquad
\Pi_{\rm col}^2=\Pi_{\rm col}=\Pi_{\rm col}^\dagger.
\end{equation*}
收集效率定义为
\begin{equation}
\boxed{
\mathcal Q_{\rm col}
\equiv\langle\psi_1|\Pi_{\rm col}|\psi_1\rangle,
\qquad
0\le\mathcal Q_{\rm col}\le1.}
\label{eq:collection-heralding-efficiency-dp19}
\end{equation}
这是一个无量纲概率，表示总单激发有多少投影到被收集模式子空间。
\end{definition}

\begin{proposition}[纯单激发的源匹配亮模与收集约化]
即使 $\Pi_{\rm col}$ 所张成的收集子空间含有很多正交光学模式，对一个\emph{固定的纯单激发态} $|\psi_1\rangle$ 而言，被收集的分量 $\Pi_{\rm col}|\psi_1\rangle$ 仍只指向这个子空间中的一个确定方向。把这一归一化方向定义成亮模 $c$，再把与它正交的全部方向并入未收集/暗子空间 $u$，便可写
\[
|\psi_1\rangle
=\frac{\mathcal A_{\rm col}^*|1_c,0_u\rangle
+\mathcal A_{\rm unc}^*|0_c,1_u\rangle}
{\sqrt{|\mathcal A_{\rm col}|^2+|\mathcal A_{\rm unc}|^2}},
\]
其中 $\mathcal A_{\rm col},\mathcal A_{\rm unc}$ 是同一条件单激发态在两组正交子空间中的复振幅，只有其比值影响归一化态。此时
\[
\mathcal Q_{\rm col}
=\frac{|\mathcal A_{\rm col}|^2}
{|\mathcal A_{\rm col}|^2+|\mathcal A_{\rm unc}|^2}.
\]
对未收集模式取偏迹后，亮模状态为
\begin{equation}
\boxed{
\rho_c
=\mathcal Q_{\rm col}|1\rangle\langle1|
+(1-\mathcal Q_{\rm col})|0\rangle\langle0|.}
\label{eq:collected_bright_mode_state_dp68}
\end{equation}
因此电子后选择可以在总光学 Hilbert 空间中预示一个单激发，而特定实验光路中的单光子概率还由 $\Pi_{\rm col}$ 的模式投影决定。



设 $\psi_1$ 是规范化纯单激发态，$\Pi_{\rm col}$ 是正交收集投影，$0<\mathcal Q_{\rm col}=\|\Pi_{\rm col}\psi_1\|^2\le1$。源匹配方向 $|1_c\rangle=\Pi_{\rm col}\psi_1/\sqrt{\mathcal Q_{\rm col}}$ 唯一，未读取的正交通道取迹后亮模态为\eqnrefs{eq:collected_bright_mode_state_dp68}的真空--单量子混合；收集空间无需一维。$\mathcal Q_{\rm col}=0$ 时只剩真空，亮方向无须定义。
\end{proposition}
\begin{proof}
令 $\Pi_{\rm col}$ 为收集子空间的正交投影，$|\psi_1\rangle$ 为归一化总单激发态，并记
\[
\mathcal Q_{\rm col}
=\langle\psi_1|\Pi_{\rm col}|\psi_1\rangle.
\]
在 $0<\mathcal Q_{\rm col}<1$ 时，分别定义归一化的收集与未收集分量
\[
|1_c\rangle
\equiv
\frac{\Pi_{\rm col}|\psi_1\rangle}{\sqrt{\mathcal Q_{\rm col}}},
\qquad
|1_u\rangle
\equiv
\frac{(\mathbb I-\Pi_{\rm col})|\psi_1\rangle}
{\sqrt{1-\mathcal Q_{\rm col}}}.
\]
由 $\Pi_{\rm col}^2=\Pi_{\rm col}$ 与
$\Pi_{\rm col}(\mathbb I-\Pi_{\rm col})=0$，立刻有
\[
\langle1_c|1_c\rangle
=\langle1_u|1_u\rangle=1,
\qquad
\langle1_c|1_u\rangle=0.
\]
因此原态具有精确正交分解
\[
|\psi_1\rangle
=\sqrt{\mathcal Q_{\rm col}}\,|1_c\rangle
+\sqrt{1-\mathcal Q_{\rm col}}\,|1_u\rangle.
\]

为了看清“唯一亮模”并不要求收集空间本身是一维，在收集子空间任选一组正交单光子基
$\{|1_\mu\rangle\}$，写
\begin{equation*}
\Pi_{\rm col}|\psi_1\rangle
=\sum_\mu c_\mu|1_\mu\rangle,
\qquad
\sum_\mu|c_\mu|^2=\mathcal Q_{\rm col}.
\end{equation*}
若 $a_\mu^\dagger|0\rangle=|1_\mu\rangle$，定义与该条件源匹配的亮模产生算符
\begin{equation*}
 c^\dagger
 \equiv
 \frac{1}{\sqrt{\mathcal Q_{\rm col}}}
 \sum_\mu c_\mu a_\mu^\dagger,
 \qquad
 [c,c^\dagger]=1.
\end{equation*}
于是
\begin{equation*}
|1_c\rangle=c^\dagger|0\rangle
=\frac{1}{\sqrt{\mathcal Q_{\rm col}}}
\sum_\mu c_\mu|1_\mu\rangle.
\end{equation*}
再把收集子空间中的其余正交线性组合补成完备模式基，它们对应的产生算符 $d_\nu^\dagger$ 都满足
$\langle0|d_\nu|\psi_1\rangle=0$。因此，多模收集空间对这个\emph{给定纯单激发态}只含一个被占据的源匹配方向；“亮模”是由态和投影共同确定的，而不是假定光路只有一个物理模式。

把 Hilbert 空间写成亮模 $c$ 与其余正交通道 $u$ 的 Fock 乘积后，上一式等价于
\[
|\psi_1\rangle
=\sqrt{\mathcal Q_{\rm col}}\,|1_c,0_u\rangle
+\ee^{\ii\varphi}
\sqrt{1-\mathcal Q_{\rm col}}\,|0_c,1_u\rangle,
\]
其中相对相位 $\varphi$ 对随后忽略 $u$ 的探测无影响。形成总密度算符并对 $u$ 偏迹，交叉项包含
$\operatorname{Tr}_u(|0_u\rangle\langle1_u|)=0$，故
\begin{align*}
\rho_c
&=\operatorname{Tr}_u|\psi_1\rangle\langle\psi_1|\\
&=\mathcal Q_{\rm col}|1_c\rangle\langle1_c|
+(1-\mathcal Q_{\rm col})|0_c\rangle\langle0_c|,
\end{align*}
即正文\eqnrefs{eq:collected_bright_mode_state_dp68}。$\mathcal Q_{\rm col}=0$ 或 $1$ 时由连续极限得到同一结果。
\end{proof}

\begin{remark}[亮模约化的纯态条件]
收集投影与其正交补给出两个正交分量。纯态在每个子空间只选出一个矢量方向，
所以完整证明中的正交分解和偏迹不要求收集空间本身只有一个模式。
若条件单激发态本身混合，各纯态分量的源匹配方向可能不同，
多模式收集态一般不能再约化为同一个纯亮模；此时应保留收集子空间上的密度算符。
\end{remark}

对于均匀折射率 $n_{\rm out}$ 的远场介质，常用数值孔径
\begin{equation*}
{\rm NA}\equiv n_{\rm out}\sin\theta_{\max}
\end{equation*}
无量纲，只规定角度接受锥；真实 $\Pi_{\rm col}$ 还必须包含偏振、频率滤波和空间模式重叠。因此相同 NA 不一定对应相同 $\mathcal Q_{\rm col}$。

参考臂把 $\rho_{01}$ 的复相位旋转到测量基底的实轴，因此近场相位进入电子计数差；环境可区分性则通过减小 $|\rho_{01}|$ 降低可见度。

约化态中的真空权重来自未收集正交模式被偏迹，而不是另行加入的探测器损耗项。


\subsection{Green 张量统一前向模型}\label{subsec:green-forward-closure-dp47}

线性介质对任意电流源的响应已经由推迟 Green 张量确定。自由电子实验只是改变了源和读出方式：电子自身作为源时，可以统计它失去的总能量，也可以只统计传到远场的辐射；外加相干光场作为源时，则读取电子获得的相干能量调制。因此三类观测共享同一个频域响应核，只在源项和探测投影上不同。

\begin{definition}[电子源与散射响应]
对给定经典电子轨迹 $\mathbf r_e(t)$，电流 Fourier 分量为
\begin{equation}
 \widetilde{\mathbf J}_e(\mathbf r,\omega)
 =q_{\mathrm e}\int\!\dd t\,
 \mathbf v(t)\delta^{(3)}[\mathbf r-\mathbf r_e(t)]\ee^{\ii\omega t}.
\label{eq:electron_current_fourier_dp47}
\end{equation}
由线性响应产生的诱导场为
\begin{equation}
\boxed{
 \widetilde{\mathbf E}_{\rm ind}(\mathbf r,\omega)
 =\ii\mu_0\omega
 \int\!\dd^3 r'\,
 \mathbf G_{\rm sc}^R(\mathbf r,\mathbf r';\omega)
 \widetilde{\mathbf J}_e(\mathbf r',\omega),}
\label{eq:induced_field_green_dp47}
\end{equation}
其中 $\mathbf G_{\rm sc}^R=\mathbf G^R-\mathbf G_0^R$ 只保留环境相对于参考真空的散射响应。

\end{definition}

\begin{proposition}[轨迹上的损失泛函]
电子单位角频率损失概率可由电子对诱导场所做功得到。这里的虚部按算符虚部理解，核为 $[\mathbf G^R(\mathbf r,\mathbf r')-\mathbf G^R(\mathbf r',\mathbf r)^\dagger]/(2\ii)$；在本章互易介质中等于逐分量虚部。对匀速直线轨迹 $\mathbf r_e(t)=\mathbf R_0+\mathbf v_0t$，写成 Green 张量双时间积分为
\begin{equation}
\boxed{
 \Gamma_{\rm EELS}(\omega)
 =\frac{\mu_0q_{\mathrm e}^2}{\pi\hbar}
 \int\!\dd t\,\dd t'\,
 \ee^{\ii\omega(t-t')}
 \mathbf v_0\!\cdot
 \operatorname{Im}\mathbf G_{\rm sc}^R
 (\mathbf r_e(t),\mathbf r_e(t');\omega)
 \cdot\mathbf v_0 .}
\label{eq:eels_green_master_dp47}
\end{equation}
$\Gamma_{\rm EELS}(\omega)\dd\omega$ 是无量纲概率，因此 $[\Gamma_{\rm EELS}]={\rm s}$。

\end{proposition}

\begin{definition}[收集通道振幅与光子谱]
CL 不是把\eqnrefs{eq:eels_green_master_dp47} 中的 $\operatorname{Im}\mathbf G^R$ 换成模方，而是先把电子电流经 Green 张量传播到远场探测通道。令探测模式 $\lambda$ 的归一化接收函数为 $\mathbf f_\lambda(\mathbf r,\omega)$，定义通道振幅
\begin{equation}
 \mathcal A_\lambda(\omega)
 =\ii\mu_0\omega
 \int\!\dd^3 r\,\dd^3 r'\,
 \mathbf f_\lambda^*(\mathbf r,\omega)\!\cdot
 \mathbf G^R(\mathbf r,\mathbf r';\omega)
 \widetilde{\mathbf J}_e(\mathbf r',\omega).
\label{eq:cl_channel_amplitude_dp47}
\end{equation}
将接收函数归一化为谱能量振幅，使 $|\mathcal A_\lambda|^2$ 的单位为 $\mathrm{J\,s}$、表示每单位角频率的通道能量后，光子谱密度为
\begin{equation}
\boxed{
 \Gamma_{\rm CL}^{\rm det}(\omega)
 =\sum_{\lambda\in\mathcal C}
 \frac{|\mathcal A_\lambda(\omega)|^2}{\hbar\omega},}
\label{eq:cl_detected_green_master_dp47}
\end{equation}
其中集合 $\mathcal C$ 编码数值孔径、偏振、频率窗和空间模式匹配。这样“总辐射”与“被探测 CL”严格区分。

\end{definition}

\begin{proposition}[相干驱动的轨迹积分]
受激 PINEM 则由外加相干场沿电子轨迹所做的功决定。对沿 $z$ 方向以速度 $v_0$ 传播的电子，时域功为
\begin{equation}
\boxed{
 \Delta E_{\rm ext}
 =q_{\mathrm e}\int_{-\infty}^{\infty}\!\dd t\,
 v_0 E_{{\rm ext},z}(\mathbf R_0,v_0t;t).}
\label{eq:pinem_time_domain_work_dp68}
\end{equation}
若外场只含角频率 $\omega$，定义正频幅 $\mathbf E_{\rm ext}^{(+)}(\mathbf r)$，使实场为 $\mathbf E_{\rm ext}^{(+)}\ee^{-\ii\omega t}+\mathrm{c.c.}$。它的单位是 ${\rm V/m}$，不同于脉冲的连续 Fourier 谱幅 ${\rm V\,s/m}$。沿轨迹的经典无量纲耦合为
\begin{equation}
\boxed{
 \beta(\omega)
 =\frac{q_{\mathrm e}}{\hbar\omega}
 \int_{-\infty}^{\infty}\!\dd z\,
 E_{{\rm ext},z}^{(+)}(\mathbf R_0,z)
 \ee^{-\ii\omega z/v_0}.}
\label{eq:pinem_beta_green_unified_dp47}
\end{equation}
若外场由具有相同正频幅约定的单频电流 $\mathbf J_{\rm ext}^{(+)}$ 经同一环境产生，则将
\begin{equation*}
 \mathbf E_{\rm ext}^{(+)}
 =\ii\mu_0\omega\,\mathbf G^R*\mathbf J_{\rm ext}^{(+)}
\end{equation*}
代入\eqnrefs{eq:pinem_beta_green_unified_dp47}，PINEM 也成为同一个 $\mathbf G^R$ 的线性泛函。若改用实场 $\operatorname{Re}[\mathbf E_{\rm ph}\ee^{-\ii\omega t}]$ 的通常相量，则 $\mathbf E_{\rm ext}^{(+)}=\mathbf E_{\rm ph}/2$；不能把连续 Fourier 电流谱直接当作该单频幅度。

\end{proposition}

\begin{remark}[同一响应核的三种读取]
因此三类实验的结构可压缩为
\begin{equation}
\boxed{
 \mathbf G^R
 \longrightarrow
 \begin{cases}
 \operatorname{Im}\mathbf G^R & \text{EELS：电子不可逆损失},\\
 \Pi_{\mathcal C}\mathbf G^R\widetilde{\mathbf J}_e & \text{CL：辐射到收集通道},\\
 \mathbf G^R\mathbf J_{\rm ext}^{(+)} & \text{PINEM：外场相干驱动},
 \end{cases}}
\label{eq:green_three_experiments_map_dp47}
\end{equation}
这里的差别来自输入源和最终读取的观测量，而不是三种不同的介质理论。

\end{remark}

\begin{definition}[仪器响应后的计数似然]
实际反演还要把理论谱映射到探测器计数。把待估参数记为 $\boldsymbol\vartheta$，理论连续谱记为 $s(\omega|\boldsymbol\vartheta)$，仪器响应矩阵为 $R_{j}(\omega)$，则第 $j$ 个像素的期望计数为
\begin{equation}
\boxed{
 \lambda_j(\boldsymbol\vartheta)
 =N_e\int\!\dd\omega\,
 R_j(\omega)s(\omega|\boldsymbol\vartheta)+b_j,}
\label{eq:detector_count_forward_dp47}
\end{equation}
其中 $N_e$ 是入射电子数，$b_j$ 是背景期望计数。若入射采用 Poisson 模型并经独立分流，像素计数独立，完整似然为
\begin{equation}
\boxed{
 \log\mathcal L(\boldsymbol\vartheta)
 =\sum_j\left[n_j\log\lambda_j-\lambda_j-\log(n_j!)\right].}
\label{eq:poisson_likelihood_green_inverse_dp47}
\end{equation}
这使 Green 张量、电子态、收集模式、仪器分辨率与最终参数反演处在同一个前向模型中。

\end{definition}

\begin{proof}[匀速损失泛函的约化]
一般的“做功 $\to\operatorname{Im}\mathbf G^R\to$ 损失概率”关系已经由\eqnrefs{eq:prescribed_current_green_loss_master,eq:green_tensor_eels_master} 建立。对具体几何，只需把相应响应核代入一般表达式：使用散射 Green 张量 $\mathbf G_{\rm sc}^R=\mathbf G^R-\mathbf G_0^R$ 去除参考真空贡献，并取匀速轨迹
\[
\mathbf r_e(t)=\mathbf R_0+\mathbf v_0t,
\qquad
\mathbf J_e(\mathbf r,t)=q_{\mathrm e}\mathbf v_0\delta^{(3)}[\mathbf r-\mathbf r_e(t)].
\]
一般表达式随即化成\eqnrefs{eq:eels_green_master_dp47}。由此可以直接看到：材料和几何只改变 Green 前向核，电子跃迁理论仍由同一个电流--响应双线性控制。
\end{proof}

\subsection{从辐射通道到有限孔径计数}
\label{chap:feqo-cl}

\begin{proposition}[偶极轴向圆锥的收集比例]
设真空远场由固定方向的电偶极 $\mathbf p=p_z\hat{\mathbf z}$ 辐射，探测器完整收集以 $+z$ 为轴、半角 $\theta_{\max}$ 的圆锥内两种偏振。其在总远场辐射中的比例为
\begin{equation}
\eta_{\rm cone}
=\frac{2-3\cos\theta_{\max}+\cos^3\theta_{\max}}4.
\label{eq:feqo-dipole-collection-fraction}
\end{equation}
\end{proposition}
\begin{proof}
偶极远场系数为
\[
\mathbf F(\hat{\mathbf n},\omega)
=\frac{k^2}{4\pi\varepsilon_0}
\hat{\mathbf n}\times(\mathbf p\times\hat{\mathbf n}),
\qquad |\mathbf F|^2\propto\sin^2\theta.
\]
因此角度功率积分给
\[
\eta_{\rm cone}
=\frac{2\pi\int_0^{\theta_{\max}}\sin^3\theta\dd\theta}
{2\pi\int_0^\pi\sin^3\theta\dd\theta},
\qquad\int_0^\pi\sin^3\theta\dd\theta=\frac43,
\]
积分后即得结论。半球给 $1/2$；窄轴向圆锥满足 $\eta_{\rm cone}=3\theta_{\max}^4/16+O(\theta_{\max}^6)$，体现偶极轴上的辐射零点。
\end{proof}

\begin{workedexample}[辐射分支与孔径效率]
若电子交出的能量中 $80\%$ 成为远场光，而探测器收集其中 $25\%$，有效收集比例为 $0.80\times0.25=0.20$。在上述轴向偶极模型中，$\theta_{\max}=60^\circ$ 给 $\eta_{\rm cone}=5/32$；若另有 $80\%$ 辐射分支，则总损失中被收集的比例为 $1/8$。

对于有导波出口的结构，稳态能量平衡应写成
\[
W_{\rm loss}=W_{\rm abs}+W_{\rm far}+W_{\rm guided},
\qquad W_{\rm det}\le W_{\rm far}.
\]
仅在无其他出口、初末局域储能相同的条件下，才能简化为吸收加远场辐射两项。数值孔径改变收集投影；材料吸收和导波损失还需由各自通道求出。
\end{workedexample}

\section{从响应核到可辨识参数}
\label{chap:feqo-end-to-end}

前面的谱与收集公式确定了理想信号。实验反演还需要明确未知量、仪器响应和计数模型。下面把这些层次逐一连接，并用可计算的核空间检验哪些参数确实能从数据确定。

\begin{definition}[前向模型与参数层次]
用 $x$ 表示待估对象，用 $\mathcal F$ 表示给定电子源、几何、频窗和探测器后的前向映射：
\begin{equation}
d=\mathcal F(x)+\epsilon.
\label{eq:feqo-forward-operator}
\end{equation}
可以选择响应核、场分布或有限维材料参数作为 $x$，但这几种选择的线性性不同。固定轨迹时，EELS 对响应核 $\operatorname{Im}\mathbf G^R$ 是线性的；响应核对介电函数和几何通常非线性。CL 需要复 Green 张量到收集通道的振幅，PINEM 需要外加相干场。
\end{definition}

\begin{proposition}[三种读出的前向泛函]
对于本章的互易介质，沿 $z$ 的匀速电子在适用的弱损失条件下满足
\begin{equation}
\Gamma_{\rm EELS}(\omega;\vb{R})
=\frac{\mu_0q_{\mathrm e}^2}{\pi\hbar}
\int\dd z\dd z'\,\ee^{\ii\omega(z-z')/v}
\operatorname{Im}G^R_{{\rm sc},zz}
((\vb{R},z),(\vb{R},z');\omega).
\label{eq:feqo-eels-forward-green}
\end{equation}
互易情形下对调 $z,z'$ 给出相同实数。固定电子源以后，核的映射为
\begin{equation}
\mathcal F_{\rm EELS}:\operatorname{Im}\mathbf G^R_{\rm sc}\longmapsto\Gamma_{\rm EELS}.
\label{eq:feqo-eels-forward-map}
\end{equation}
对于具有通常真空远场展开的有限结构，定义
\[
\mathbf G_\infty(\hat{\vb{n}},\vb{r}';\omega)
=\lim_{r\to\infty}r\ee^{-\ii kr}\mathbf G^R(r\hat{\vb{n}},\vb{r}';\omega),
\quad
\vb{F}=\ii\mu_0\omega\int\mathbf G_\infty\widetilde{\vb{J}}_e\dd^3r'.
\]
取 $\vb{E}(t)=(2\pi)^{-1}\int\widetilde{\vb{E}}(\omega)\ee^{-\ii\omega t}\dd\omega$。若 $\mathsf Q(\hat{\vb{n}},\omega)$ 是偏振空间中满足 $0\le\mathsf Q\le I$ 的接收效率矩阵，则
\begin{equation}
\Gamma_{\rm CL}^{\rm det}
=\frac{1}{\pi Z_0\hbar\omega}
\int\dd\Omega\,\vb{F}^\dagger\mathsf Q\vb{F},
\qquad Z_0=\sqrt{\mu_0/\varepsilon_0}.
\label{eq:feqo-cl-forward-collection}
\end{equation}
由此
\begin{equation}
\mathcal F_{\rm CL}:\mathbf G_\infty\longmapsto\Gamma_{\rm CL}^{\rm det}.
\label{eq:feqo-cl-forward-map}
\end{equation}
只给出 $\operatorname{Im}G_{{\rm sc},zz}$ 不能确定这一通道振幅。对外加场，沿用本书正频复幅定义，有
\begin{align}
\beta(\vb{R})&=\frac{q_{\mathrm e}}{\hbar\omega}
\int E^{(+)}_z(\vb{R},z)\ee^{-\ii\omega z/v}\dd z,
\label{eq:feqo-pinem-forward-beta}\\
\mathcal F_{\rm PINEM}:E^{(+)}_z&\longmapsto
\{J_\ell^2(2|\beta|)\}_{\ell\in\ZZ}.
\label{eq:feqo-pinem-forward-map}
\end{align}
最后一式使用窄带、无反冲能量格和经典相干驱动条件。
\end{proposition}
\begin{proof}
第一式由电子电流代入\eqnrefs{eq:eels_green_master_dp47}，使用 $\dd z=v\dd t$ 得到。远场功率为 $r^2|\vb{E}|^2/Z_0$，正频 Parseval 公式给系数 $1/\pi$；除以单光子能量 $\hbar\omega$ 并投影接收偏振，得到 CL 式。PINEM 的线积分与边带概率则分别由\eqnrefs{eq:pinem_beta_green_unified_dp47}及经典驱动解得到。
\end{proof}

\begin{workedexample}[有限窗口核的伴随]
固定一个频率，将轨迹截在 $[-L/2,L/2]$，未知核取 $H\in L^2([-L/2,L/2]^2)$。记 $C=\mu_0q_{\mathrm e}^2/(\pi\hbar)$，则
\[
TH=C\int\dd z\dd z'\,\ee^{\ii\omega(z-z')/v}H(z,z'),
\qquad |TH|\le |C|L\|H\|_2.
\]
求该前向算子的伴随，并说明有限轨迹模型为何给出有界映射。

\textbf{解。} 取内积 $\langle f,g\rangle=\int f^*g$，有
\[
(T^*y)(z,z')=C^*\ee^{-\ii\omega(z-z')/v}y.
\]
这才是一个有界算子的伴随。在整个六维体积 $L^2$ 空间中，“把核限制到理想零宽轨迹”并不自动有界；实际束宽、空间正则性或有限维离散化必须进入模型。多位置观测对应多个这样的轨迹核，其伴随为各核的加权和。
\end{workedexample}

\subsection{参考相位、仪器响应与计数}

\begin{workedexample}[高斯近场的三相位复振幅读取]
设某条轨迹上的正频场为
\[
E_z^{(+)}(z)=E_*\ee^{\ii\phi}\ee^{-z^2/(2d^2)}\ee^{\ii k_z z},
\qquad\Delta k=k_z-\omega/v.
\]
求该轨迹上的复耦合，并用已标定参考的三个相位恢复它的两个正交分量；比较弱耦合计数与精确边带计数。

\textbf{解。} 直接完成高斯积分，得到
\begin{equation}
\beta_s=\frac{q_{\mathrm e}E_*\sqrt{2\pi}d}{\hbar\omega}
\ee^{-d^2\Delta k^2/2}\ee^{\ii\phi}
\equiv A\ee^{\ii\varphi}.
\label{eq:feqo-imaging-beta}
\end{equation}
$A\ge0$；电子电荷 $q_{\mathrm e}=-e$ 的负号并入 $\varphi$。沿同一相干电子轨迹叠加已标定参考耦合 $B\ee^{\ii\theta}$，则 $\beta_\theta=\beta_s+B\ee^{\ii\theta}$。记 $r_\theta=|\beta_\theta|$，弱耦合时
\begin{equation}
P_{+1}=P_{-1}=r_\theta^2+O(r_\theta^4),
\qquad P_0=1-2r_\theta^2+O(r_\theta^4).
\label{eq:feqo-imaging-weak-forward}
\end{equation}
写 $\beta_s=x+\ii y$，理想弱耦合强度 $Y_\theta=r_\theta^2$ 满足
\[
Y_\theta=x^2+y^2+B^2+2B(x\cos\theta+y\sin\theta).
\]
对 $\theta=0,\pi/2,\pi$，令 $C_0=(Y_0+Y_\pi)/2$，便有
\[
x=\frac{Y_0-Y_\pi}{4B},\qquad
y=\frac{Y_{\pi/2}-C_0}{2B},\qquad B>0.
\]
取 $A=B=0.05$、$\varphi=\pi/3$，三次测量的 $Y_\theta$ 为
\[
0.0075,\qquad0.005+0.0025\sqrt3\simeq0.00933013,
\qquad0.0025.
\]
每个相位入射 $10^5$ 个电子时，单侧一级边带的弱耦合期望计数约为 $750,933,250$。精确拟合应使用 $J_1^2(2r_\theta)$；相应计数约为 $744,924,249$，体现截断误差已经可见。
\end{workedexample}

\begin{definition}[归一化探测响应与计数模型]
若一级边带间距为 $\hbar\omega$，第 $\ell$ 条真实谱线经高斯分辨函数落入像素窗 $\mathcal B_j$ 的概率为
\[
R_{j\ell}=\int_{\mathcal B_j}\frac{\dd E}{\sqrt{2\pi}\sigma_E}
\exp\!\left[-\frac{(E-E_0-\ell\hbar\omega)^2}{2\sigma_E^2}\right].
\]
加入未检测结果后，每列满足 $\sum_jR_{j\ell}=1$。相位 $\theta_s$ 下的像素概率与 Poisson 入射模型为
\begin{equation}
p_{js}=\sum_\ell R_{j\ell}J_\ell^2(2|\beta_{\theta_s}|),
\qquad\mu_{js}=\nu_s p_{js}+b_{js},\qquad
n_{js}\sim\operatorname{Poisson}(\mu_{js}).
\label{eq:feqo-imaging-detector}
\end{equation}
独立 Poisson 计数适用于独立 Poisson 入射经分流的模型。若入射总数固定，各出口计数服从多项分布，有负协方差；不能把固定总数数据直接当作独立 Poisson 样本。
\end{definition}

\begin{proposition}[局部可辨识性的 Fisher 判据]
设有限维实参数 $\vb{\vartheta}$ 的所有 $\mu_j>0$，导数存在，且采用独立 Poisson 模型。令 $J_{ja}=\partial_a\mu_j$，则
\begin{equation}
I=J^{\mathsf T}\operatorname{diag}(\mu_j^{-1})J,
\qquad\operatorname{Cov}(\widehat{\vb{\vartheta}})\succeq I^{-1}.
\label{eq:feqo-cramer-rao}
\end{equation}
后一不等式要求 $I$ 可逆、估计无偏且满足通常微分积分互换条件。$I$ 正定当且仅当加权雅可比矩阵满列秩。
\end{proposition}
\begin{proof}
分数函数为 $\partial_a\log\mathcal L=\sum_j(n_j/\mu_j-1)\partial_a\mu_j$。利用 $\operatorname{Var}n_j=\mu_j$ 和独立性得到 Gram 矩阵 $I$。对任意向量 $h$，
\[
h^{\mathsf T}Ih=\sum_j\frac{(J_jh)^2}{\mu_j}.
\]
因此存在不可辨识的一阶方向当且仅当 $Jh=0$。Cram\'er--Rao 界由估计误差与分数函数的协方差恒等式及矩阵 Cauchy--Schwarz 不等式得到。
\end{proof}

\begin{figure}[htbp]
\centering
\includegraphics[width=.82\linewidth]{electron_field_calibration.pdf}
\caption{相位失配压低有效耦合，完整 Bessel 分布决定可见边带。高阶边带显著时，弱耦合近似不能用于定量反演；场幅估计还需校准轨迹积分与仪器响应。}
\label{fig:electron-field-calibration-modern}
\end{figure}

\subsection{正则化、噪声放大与空间分辨}

\begin{theorem}[有限维 Tikhonov 解与唯一性]
对白化并统一量纲后的数据 $d\in\CC^m$、未知系数 $x\in\CC^n$，给定 $A,L$ 和 $\eta>0$，定义
\begin{equation}
\widehat x_\eta=\mathop{\rm argmin}_x
\{\|Ax-d\|_2^2+\eta\|Lx\|_2^2\}.
\label{eq:feqo-tikhonov-regularization}
\end{equation}
解唯一当且仅当 $\ker A\cap\ker L=\{0\}$；此时
\begin{equation}
(A^\dagger A+\eta L^\dagger L)\widehat x_\eta=A^\dagger d.
\label{eq:feqo-tikhonov-normal}
\end{equation}
\end{theorem}
\begin{proof}
对实部和虚部求变分得到正规方程。二次型满足
\[
x^\dagger(A^\dagger A+\eta L^\dagger L)x
=\|Ax\|_2^2+\eta\|Lx\|_2^2,
\]
故正定性等价于公共核仅含零向量；若存在公共零方向，目标函数沿该方向不变。
\end{proof}

\begin{proposition}[奇异值滤波与误差界]
取 $L=I$，$A$ 的秩为 $r$，非零奇异值与奇异向量为 $\sigma_j,u_j,v_j$，$1\le j\le r$。真实数据 $d=Ax_*+\epsilon$、$\|\epsilon\|\le\delta$，$P_{\ker A}$ 为未知空间中核的正交投影。则
\begin{equation}
\widehat x_\eta=\sum_{j=1}^r\frac{\sigma_j}{\sigma_j^2+\eta}
\langle u_j,d\rangle v_j,
\label{eq:feqo-imaging-field-inversion}
\end{equation}
并且
\begin{align}
\|\widehat x_\eta-x_*\|
&\le\frac{\delta}{2\sqrt\eta}+\|b_\eta\|,\label{eq:feqo-tikhonov-error-bound}\\
b_\eta&=P_{\ker A}x_*+
\sum_{j=1}^r\frac{\eta}{\sigma_j^2+\eta}\langle v_j,x_*\rangle v_j.
\label{eq:feqo-tikhonov-bias}
\end{align}
第一项偏差 $P_{\ker A}x_*$ 完全不可从数据恢复。若额外有源条件 $x_*=A^\dagger z$，该项为零，总偏差不超过 $\sqrt\eta\|z\|/2$。
\end{proposition}
\begin{proof}
逐个奇异向量求正规方程得到滤波式。标量函数 $\sigma/(\sigma^2+\eta)$ 的最大值为 $1/(2\sqrt\eta)$，控制噪声放大；同一函数乘 $\eta$ 控制源条件下的偏差。没有源条件时，偏差必须保留其显式谱分解。
\end{proof}

\begin{workedexample}[高斯模糊与边带信息的秩]
设仪器空间响应为宽度 $w$ 的归一化高斯核，未知场离散为 $n$ 个复系数。求噪声为 $\delta_k$ 时可分辨的空间频率条件，以及固定轨迹、固定频率 PINEM 计数能够约束的实参数数目。

\textbf{解。} 一维高斯空间模糊的 Fourier 奇异值为 $\sigma(k)=\ee^{-w^2k^2/2}$。某一空间 Fourier 分量只有在 $\sigma(k)|x_k|\gtrsim\delta_k$ 时可分辨，因此信噪比大于一时满足
\[
|k|\lesssim\frac1w\sqrt{2\log(|x_k|/\delta_k)}.
\]
这个判据同时依赖仪器宽度和该分量的信噪比。

\begin{samepage}
在固定轨迹、固定频率下，任意多个 PINEM 边带的概率都只通过一个复数 $\beta[E]$ 依赖场。若离散场有 $n$ 个复系数，计数对其 $2n$ 个实坐标的雅可比矩阵必须因式分解为
\[
J_{\rm counts,field}=J_{\rm counts,\beta}J_{\beta,\rm field},
\qquad\operatorname{rank}J_{\rm counts,field}\le2.
\]
\begingroup\interlinepenalty=10000
无参考时还只依赖 $|\beta|$，局部秩至多一。高阶边带改善同一耦合的统计估计；增加独立空间核需改变轨迹、频率、入射场或观测通道。\par
\endgroup
\end{samepage}
\end{workedexample}

\subsection{球体近场的完整损失谱与参数退化}

平面例题将响应分成材料因子与轨迹因子。小球的偶极极限也有同样分解，可用来检验“改变扫描距离就能同时测出几何与介电函数”的说法。关于电子损失与纳米颗粒响应的完整多极理论，参见\cite{GarciaDeAbajo2010RMP}。

\begin{proposition}[非迟滞偶极球的 EELS 谱]
设半径 $a$ 的均匀非磁性球位于原点，电子沿 $(b,0,vt)$ 运动。假定 $a/b\ll1$、$\omega a/c\ll1$、$|\sqrt\varepsilon|\omega a/c\ll1$、$v/c\ll1$，并采用弱损失线性响应。材料侧的条件控制内部场的准静态约化，不能从真空小尺寸条件直接推出。定义体积量纲极化率
\[
\alpha_V=4\pi a^3\frac{\varepsilon-1}{\varepsilon+2},
\qquad\vb{p}=\varepsilon_0\alpha_V\vb{E}_e,
\qquad s=\omega b/v.
\]
则单电子、单位角频率的偶极损失概率为
\begin{equation}
{\Gamma_{\rm sph}(\omega,b)
=\frac{q_{\mathrm e}^2\omega^2}{4\pi^3\varepsilon_0\hbar v^4}
\operatorname{Im}\alpha_V(\omega)
[K_0^2(s)+K_1^2(s)].}
\label{eq:feqo-sphere-eels-dipole}
\end{equation}
若保留辐射损失到对应精度，极化率还需包含辐射反作用修正；在窄共振邻域，动态去极化和辐射修正还须与小分母 $\varepsilon+2$ 比较。本式的非迟滞近似不宜直接用于高速电子或 $b$ 接近 $a$ 的高多极贡献。
\end{proposition}
\begin{proof}
原点的电子库仑场为
\[
E_{e,x}(t)=-\frac{q_{\mathrm e}b}{4\pi\varepsilon_0(b^2+v^2t^2)^{3/2}},
\qquad
E_{e,z}(t)=-\frac{q_{\mathrm e}vt}{4\pi\varepsilon_0(b^2+v^2t^2)^{3/2}}.
\]
从
\[
\int_{-\infty}^{\infty}\frac{\ee^{\ii\omega t}\dd t}{\sqrt{b^2+v^2t^2}}
=\frac2vK_0(s)
\]
分别对 $b$ 求导、对 $t$ 分部积分，得到
\[
\widetilde E_{e,x}=-\frac{q_{\mathrm e}\omega}{2\pi\varepsilon_0v^2}K_1(s),
\qquad
\widetilde E_{e,z}=-\frac{\ii q_{\mathrm e}\omega}{2\pi\varepsilon_0v^2}K_0(s).
\]
偶极吸收的总功为 $W=\int\vb{E}_e\cdot\dot{\vb{p}}\dd t$。实场的正频 Parseval 恒等式给
\[
\frac{\dd W}{\dd\omega}
=\frac{\varepsilon_0\omega}{\pi}
\operatorname{Im}\alpha_V|\widetilde{\vb{E}}_e|^2.
\]
代入两场分量，再除以 $\hbar\omega$ 得到结论。这里两个轨迹因子来自激发场和响应场的双线性，不能用单一源点的场幅替代。
\end{proof}

\begin{workedexample}[距离标度、Drude 球峰与半径估计]
设纳米球满足前述电偶极近似，并采用局域 Drude 介电函数。求小、大冲击参数下的损失标度和弱损耗材料极点；介电函数已独立标定时，由极化率的虚部反演球半径。

\textbf{解。} 当 $s\ll1$ 时，$K_1(s)\sim1/s$、$K_0(s)\sim-\log(s/2)-\gamma_{\rm E}$，故
\[
\Gamma_{\rm sph}\sim\frac{q_{\mathrm e}^2}{4\pi^3\varepsilon_0\hbar v^2b^2}
\operatorname{Im}\alpha_V.
\]
距离主标度为 $b^{-2}$。当 $s\gg1$ 时，$K_0^2+K_1^2\sim(\pi/s)\ee^{-2s}$，损失具有指数距离截止。

对 Drude 球，令 $f=(\varepsilon-1)/(\varepsilon+2)$，直接代入得到
\[
f=-\frac{\omega_p^2}{3\omega^2+3\ii\Gamma_D\omega-\omega_p^2},
\qquad
\operatorname{Im}f=
\frac{3\Gamma_D\omega\omega_p^2}
{(3\omega^2-\omega_p^2)^2+9\Gamma_D^2\omega^2}.
\]
弱损耗的材料极点在 $\omega_p/\sqrt3$；测得谱峰还受 Bessel 轨迹权重影响。若介电函数已经独立标定，测得 $\operatorname{Im}\alpha_V$ 后可用
\[
a=\left[\frac{\operatorname{Im}\alpha_V}{4\pi\operatorname{Im}f}\right]^{1/3}
\]
恢复半径。在 $\operatorname{Im}f$ 很小的频段，这一估计对噪声和材料标定误差敏感。
\end{workedexample}

\begin{proposition}[固定频率的半径--材料退化]
在同一频率下，未知半径 $a$ 与未知材料因子 $S=\operatorname{Im}f>0$ 只以乘积 $a^3S$ 出现。任意多个距离 $b_j$ 的独立弱损失计数模型 $\mu_j=C_j a^3S$，其关于 $(a,S)$ 的 Fisher 矩阵秩至多为一。
\end{proposition}
\begin{proof}
有 $\partial_a\mu_j=3\mu_j/a$、$\partial_S\mu_j=\mu_j/S$，故雅可比矩阵的两列比例相同，方向
\[
\delta S=-3S\delta a/a
\]
使所有一阶计数变化为零。事实上变换 $a\mapsto ca$、$S\mapsto S/c^3$ 保持全部理想距离扫描数据不变。

因此距离扫描可检验轨迹因子与偶极近似，并估计乘积幅度；分离半径和任意介电函数还需独立几何标定、受约束材料模型、多极响应或额外观测。频率扫描是否解除退化，必须检查具体参数化模型的雅可比矩阵。
\end{proof}


\begin{derivation}[外场做功与近场复耦合]
采用全文的正频率复幅约定，把真实外场写成
\[
E_{{\rm ext},z}(\vb{r},t)
=\int_0^\infty\frac{\dd\omega}{2\pi}
\left[
E_{{\rm ext},\omega,z}(\vb{r})\ee^{-\ii\omega t}
+E_{{\rm ext},\omega,z}^*(\vb{r})\ee^{+\ii\omega t}
\right].
\]
沿匀速轨迹 $\vb{r}_e(t)=(\vb{R}_0,v_0t)$ 代入正文\eqnrefs{eq:pinem_time_domain_work_dp68}。对固定正频率 $\omega$，定义复功幅
\begin{align*}
\mathcal W_\omega
&\equiv
q_{\mathrm e}\int_{-\infty}^{\infty}\dd t\,
 v_0 E_{{\rm ext},\omega,z}(\vb{R}_0,v_0t)
 \ee^{-\ii\omega t}\\
&=q_{\mathrm e}\int_{-\infty}^{\infty}\dd z\,
 E_{{\rm ext},\omega,z}(\vb{R}_0,z)
 \ee^{-\ii\omega z/v_0},
\end{align*}
其中第二行只用了 $z=v_0t$ 与 $\dd z=v_0\dd t$。指数因子表明电子只抽取外场纵向 Fourier 谱在
\[
k_z=\frac{\omega}{v_0}
\]
这一同步线上对应的分量。

单个边带步长的能量是 $\hbar\omega$，因此把相干功幅除以这一能量定义无量纲生成元
\[
\beta(\omega)\equiv\frac{\mathcal W_\omega}{\hbar\omega},
\]
立刻得到正文\eqnrefs{eq:pinem_beta_green_unified_dp47}。这个除法不是经验归一化：在能量格演化算符中，$\beta$ 必须是乘在无量纲平移生成元前的无量纲系数。

若外场由外源电流产生，正文采用
\[
\vb{E}_{{\rm ext},\omega}(\vb{r})
=\ii\mu_0\omega
\int\dd^3r'\,
\mathbf G^R(\vb{r},\vb{r}';\omega)
\cdot\widetilde{\vb{J}}_{\rm ext}(\vb{r}',\omega).
\]
把它代回 $\beta$ 并约去 $\omega$，得到
\begin{align*}
\beta(\omega)
={}&\frac{\ii\mu_0q_{\mathrm e}}{\hbar}
\int\dd z\dd^3r'\,
 \ee^{-\ii\omega z/v_0}
\hat{\vb{z}}\cdot
\mathbf G^R((\vb{R}_0,z),\vb{r}';\omega)
\cdot\widetilde{\vb{J}}_{\rm ext}(\vb{r}',\omega).
\end{align*}
因此 PINEM 耦合与 EELS/CL 使用的是同一个推迟 Green 算符；PINEM 由外源产生相干场并沿电子轨迹读取线性振幅，EELS 则对电子自身诱导场取耗散二次型。
\end{derivation}

